title: Unit Tests
alias: unit tests
todo:
------------------------------
\usepackage{listings}
\lstset{language=C++,basicstyle=\small\ttfamily,backgroundcolor=\color{black!10!white},columns=fullflexible,tabsize=4}
------------------------------

GDSiMS uses the \href{https://github.com/catchorg/Catch2}{Catch2} testing framework (version 3), fetched automatically by CMake the first time the tests are built. The entire suite is built and run with a single [[make options|command]]:

\begin{Code}
	make test
\end{Code}

The tests are built in debug mode, with the sanitizers on, so a test which (for example) reads past the end of an array fails loudly, even if it would have happened to give the right answer.

\section{Layout}

	The test sources live in _tests/_, one file per area of the code. Every _.cpp_ file in that directory is picked up automatically, so a new test file needs no changes to the build.

	The tests link against the same static library as the executable, so they test the real code, rather than a copy of it. Code in _src/app/_ is not part of that library: anything worth testing should therefore live in the library, with _src/app/_ kept as thin as possible.

	The compiled test binary is placed at _tests/gdsims\_tests_, and can be run directly, which gives access to Catch2's own options; for example, to run only the tests with a given tag:

	\begin{Code}
		./tests/gdsims\_tests "[validation]"
	\end{Code}

\section{Writing tests}

	A test checks one behaviour, and is named for that behaviour:

	\begin{lstlisting}
TEST_CASE("A thread count of zero is an error, but not fatal", "[validation]")
{
	MainSettings s;
	s.Runtime.ParallelThreads = 0;
	auto status = querySettingsConstraints(s);
	REQUIRE(status.Errors.size() == 1);
	REQUIRE(status.Fatal.empty());
}
\end{lstlisting}

	Some things to bear in mind:

	* Tests should construct their own settings, rather than modifying the global _Settings_ object, so that they cannot affect one another.
	* Behaviour which depends on the machine (the number of cores, whether a terminal is attached) should either be skipped (with _SKIP_) when the precondition does not hold, or not tested at all.
	* Anything stochastic should be tested with a fixed seed, or with a statistical tolerance wide enough that the test does not fail by chance.

